% Page 037

\seite{37}

die Schule besichtigt und seine Zufrieden\ohb
heit ausgesprochen wurde.

\pend
\abschnitt{Wechsel in der Ortsschulinspection}
\pstart

Pfarrer Ettges wurde am 14.\ Mai durch\ob
Verfügung der Königlichen Regierung abgelöst.\ob
Herr Pfarrer Schmitt zu Seffern wurde zum\ob
Ortsschulinspector ernannt.

\pend
\abschnitt{Ernte}
\pstart

Durch die oft sehr lange anhaltende regnerische\ob
Witterung wurde die diesjährige Ernte sehr\ob
weit ausgedehnt. Die Ernte hat ununterbrochen\ob
stattgefunden. Verdorben fast der größte Teil\ob
der Obstfrüchte, und infolge dessen wuchs\ob
nur sehr wenig Ohmd. Während in früheren\ob
Jahren die Kartoffeln meist im April gesetzt\ob
wurden, mußte diesmal des Wetters wegen\ob
die Kartoffelsaat Ende Mai vollzogen werden.\ob
Das Heu machte, da es während der Heuernte\ob
öfter regnete, vielfach nur als minder\ohb
wertig eingeheimst. Roggen und Hafer blieben\ob
teilweise bis spät in den Herbst hinein auf dem\ob
Felde stehen. Sehr viel Futter, Hafer und Roggen\ob
mußten in den Garben auf den Wagen liegen,\ob
anstatt an trockener Stelle auf- und abgeladen\ob
zu werden. Die Folge war, daß die Früchte\ob
in verfaultem Zustande auf dem Felde liegen\ob
blieben. Am Tage Allerheiligen war noch eine\ob
große Menge Kartoffeln auszufahren. Die Preise\ob
waren nur in geringer Quantität gewachsen,\ob
und der Preis derselben stieg bis 3 M für 50 kg.\ob
Ungemein tief standen, im Gegensatze zu den\ob
Vorjahren, die Getreidepreise.

\pend
\abschnitt{1/}
\pstart


\pend
\abschnitt{Schulfeierlichkeiten}
\pstart

Der Geburtstag Seiner Majestät des Kaisers\ob
wurde von den vier Schulen der Pfarrei\ob
Seffern, unter Beisein des Herrn Ortsschul\ohb
inspectors, Herrn Pfarrer Schmitt, und sämt\ohb
licher Lehrpersonen gemeinsam in Seffern\ob
begangen. Abwechselnd mit Liedern, die
\pend
